% Generated draft: requires mathematical and source review.
\documentclass[11pt]{article}
\usepackage[margin=25mm]{geometry}
\usepackage{fontspec}
\setmainfont{Latin Modern Roman}
\usepackage{amsmath,amssymb,mathtools,unicode-math}
\setmathfont{Latin Modern Math}
\usepackage{enumitem,needspace,xurl,hyperref}
\setlength{\parindent}{0pt}
\setlength{\parskip}{5pt}
\setlength{\emergencystretch}{3em}
\newcommand{\N}{\mathbb{N}}
\newcommand{\Z}{\mathbb{Z}}
\newcommand{\Q}{\mathbb{Q}}
\newcommand{\R}{\mathbb{R}}
\newcommand{\C}{\mathbb{C}}
\newcommand{\F}{\mathbb{F}}
\newcommand{\A}{\mathbb{A}}
\newcommand{\PP}{\mathbb{P}}
\newcommand{\E}{\mathbb{E}}
\newcommand{\eps}{\varepsilon}
\newcommand{\dd}{\,\mathrm{d}}
\begin{document}
\section*{UnsolvedMath ID: 700922}
\textbf{Source problem code:} catalogue entry
\subsection{Definitions and mathematical statement}
For a positive integer \(n\), let \(\mathrm{GL}_n(2)=\mathrm{GL}_n(\F_2)\), acting naturally on the column-vector space \(\F_2^n\). The source considers a group \(G\) with a normal ``(pro-)\(2\)-subgroup'' \(N\). In the finite \(2\)-group interpretation, the Frattini subgroup \(\Phi(N)\) is the intersection of the maximal proper subgroups of \(N\). In the pro-\(2\) interpretation, a pro-\(2\) group is an inverse limit of finite \(2\)-groups, and \(\Phi(N)\) is the intersection of its maximal proper open subgroups. The source does not fully specify the ambient topological category for \(G\).

The hypotheses are that there exist isomorphisms
\[
\alpha:G/N\longrightarrow\mathrm{GL}_n(2),\qquad
\beta:N/\Phi(N)\longrightarrow(\F_2^n,+)
\]
such that the action induced by conjugation is the natural module action, explicitly
\[
\beta\bigl(gxg^{-1}\Phi(N)\bigr)
=\alpha(gN)\,\beta\bigl(x\Phi(N)\bigr)
\qquad(g\in G,\ x\in N).
\]
Here the action factors through \(G/N\), since inner automorphisms of \(N\) act trivially on its Frattini quotient in the finite and pro-\(2\) settings.

\textbf{Problem (B. Baumann, 1995).} For \(n=3\), determine \(G\) satisfying these hypotheses such that \(N\) is ``as large as possible.'' The source does not define this comparison: in particular, it does not specify whether to maximize finite order, seek a universal pro-\(2\) extension, or use another extremal criterion. Consequently, the extremal target remains unresolved at the level of formulation.

The source supplies the following contextual assertions: for \(n>3\), it can be proved that \(\Phi(N)=1\); for \(n=2\), there exists such a group \(G\) with \(N\) the free pro-\(2\) group on two generators. These are retained as source assertions, without independent verification. A free pro-\(2\) group on two generators has the universal mapping property for assignments of its two generators to arbitrary pro-\(2\) groups, with continuous homomorphisms.
\subsection{Short English statement}
Determine the groups \(G\) whose normal \((\mathrm{pro}\text{-})2\)-subgroup \(N\) has quotient \(G/N\cong\mathrm{GL}_3(2)\), acting naturally on \(N/\Phi(N)\), with \(N\) as large as possible. The meaning of ``as large as possible'' and the precise finite or topological category require clarification. The supplied registry labels the problem open; this is not an independently verified status assessment.
\subsection{Sources}
\begin{itemize}
\item Local UnsolvedMath registry, permanent ID \(700922\), source code \texttt{catalogue entry}, title \emph{Kourovka Notebook Problem 13.4}. The supplied record labels the problem open and gives upstream revision \texttt{d371ee917ac9647fbb06cd6fc529deec6f3dda82}.
\item \emph{The Kourovka Notebook (Unsolved Problems in Group Theory)}, original problem 13.4, proposed by B. Baumann in the 13th issue (1995). The supplied background also identifies the 21st-issue edition (2026), arXiv identifier \texttt{1401.0300v48}.
\item The supplied collection description attributes editing of \emph{Kourovka Notebook}, Issues 1--21 (1965--2026), to E. I. Khukhro and V. D. Mazurov. These bibliographic attributions and the mathematical contextual assertions are supplied source information; the cited edition has not been independently consulted here.
\end{itemize}
\end{document}
